\section*{Chapitre \ref{ch:nw:network-cards-and-kernel-bypass} --- Cartes réseau et contournement du noyau}

\begin{solution}{exo:nw:network-cards-and-kernel-bypass:1}
$1.5 \times 10^6 \times 400 \times 10^{-6} = 600$ arrivées pour $512 - 40 = 472$ descripteurs libres~: 128 paquets perdus.
\end{solution}

\begin{solution}{exo:nw:network-cards-and-kernel-bypass:2}
$1 + \lambda u = 1 + 0.2 \times 20 = 5$ paquets par interruption~; interruptions à $200\,000 / 5 = 40\,000$ par seconde.
\end{solution}

\begin{solution}{exo:nw:network-cards-and-kernel-bypass:3}
$15.4 / 16.2$~: 95~\%. Presque tout est dû au réveil et à l'ordonnancement du lecteur endormi, et non à la pile réseau.
\end{solution}

\begin{solution}{exo:nw:network-cards-and-kernel-bypass:4}
La \omterm{def:nw:network-cards-and-kernel-bypass:rss}{répartition côté réception} (RSS) répartit les flux sur les anneaux pour équilibrer les cœurs~; un feed handler veut que tous les groupes d'une ligne
tiennent sur l'unique anneau que son cœur épinglé scrute, sans autre trafic. L'\omterm{def:nw:network-cards-and-kernel-bypass:rss}{aiguillage des flux} (flow steering) affecte les flux de la plateforme à l'anneau choisi et répartit le reste du trafic de l'hôte sur les autres.
\end{solution}

\begin{solution}{exo:nw:network-cards-and-kernel-bypass:5}
$1/(2 + 0.3) = 0.435$ paquet par microseconde~: 435\,000 par seconde. Avec 16 paquets par interruption, $1/(2/16 + 0.3) = 2.35$ millions par seconde.
\end{solution}

\begin{solution}{exo:nw:network-cards-and-kernel-bypass:6}
Le fil de scrutation ne s'endort jamais, donc le réveil (l'essentiel de la médiane de blocage) disparaît. Mais sur des cœurs partagés, l'ordonnanceur exécute parfois
autre chose sur le cœur du fil de scrutation pendant des millisecondes, et chaque datagramme arrivant entre-temps attend. Sur un hôte de production, le fil de scrutation
reçoit un cœur isolé, avec ses interruptions et son tick de minuterie déportés (livre~13, chapitre~13), sur le socket de la carte.
\end{solution}

\begin{solution}{exo:nw:network-cards-and-kernel-bypass:7}
128 descripteurs~: avec des lots de 32, un anneau de 64 perd encore 36 paquets du jeu de test, un anneau de 128 n'en perd aucun.
\end{solution}

\begin{solution}{exo:nw:network-cards-and-kernel-bypass:8}
Les pertes sont rares, sans être impossibles (un module optique instable, un tampon du commutateur pendant une rafale, une bascule (failover))~: sans retransmission, un seul segment perdu
laisse la session bloquée ou manque silencieusement un ordre, et la couche session de la plateforme déconnectera ou annulera. Les deux microsecondes
viennent du retrait du noyau, non du retrait de la fiabilité~; il faut recourir à un TCP matériel ou à une \omterm{def:nw:network-cards-and-kernel-bypass:pmd}{pile réseau en espace utilisateur} qui l'implémente.
\end{solution}

\begin{solution}{pb:nw:network-cards-and-kernel-bypass:1}
\begin{enumerate}
\item $0.2 \times (2 + 0.3) = 0.46$~: 46~\% d'un cœur.
\item $3 + 2 + 0.3 = \qty{5.3}{\micro\second}$.
\item File d'attente~: une interruption qui arrive alors que la précédente est encore traitée attend que celle-ci soit terminée.
\item $1/2.3$ par microseconde~: 435\,000 paquets par seconde.
\item $1 + 0.2 \times 10 = 3$.
\item $200\,000 / 3 \approx 66\,667$ par seconde.
\item $66\,667 \times 2 \times 10^{-6} + 0.2 \times 0.3 = 0.133 + 0.06$~: 19.3~\%.
\item Avec la minuterie~: le premier paquet d'un lot attend \qty{10}{\micro\second}, les suivants moins, avant le traitement
  de \qty{5.3}{\micro\second}.
\item Une demi-boucle, \qty{0.1}{\micro\second}.
\item $0.1 + 0.3 = \qty{0.4}{\micro\second}$~; le 0.41 du modèle ajoute le rare paquet qui arrive pendant le traitement du précédent.
\item Le cœur entier.
\item Isolé, jamais désordonnancé, sur le socket de la carte, sans rien d'autre à faire.
\item \textbf{Résultat nommé.} À 200\,000 paquets par seconde, une minuterie de \qty{10}{\micro\second} donne une latence moyenne de \qty{12.4}{\micro\second}
  pour 19.3~\% d'un cœur~; la scrutation donne \qty{0.41}{\micro\second} pour la totalité~; les interruptions seules prendraient un cœur entier à 435\,000
  paquets par seconde.
\item Environ \qty{12}{\micro\second} (12.4 contre 0.41).
\item Au point où $\lambda(2/(1 + 10\lambda) + 0.3) = 0.25$~: environ 324\,000 paquets par seconde.
\item Le chemin critique (hot path) mérite un cœur~; le reste du trafic de l'hôte ne le mérite pas, et le \omterm{def:nw:network-cards-and-kernel-bypass:coalesce}{regroupement des interruptions} l'empêche d'interrompre la machine toute la journée.
\item Coût des interruptions et du réveil avec un profileur et des horodatages matériels, le comportement de regroupement de la carte sous le flux réel, et la
  distribution des pauses du cœur de scrutation.
\item $1\,024 / 200\,000$ secondes~: \qty{5.1}{\milli\second} de pause avant le débordement de l'anneau.
\item Des microsecondes de réveil par lot, et une queue de distribution aussi longue que la plus longue décision de l'ordonnanceur.
\item Accès brut à l'anneau pour le flux multicast, scruté sur un cœur isolé~: le flux n'a pas besoin de TCP, et l'accès direct en supprime le plus.
\end{enumerate}
\end{solution}

\begin{solution}{iq:nw:network-cards-and-kernel-bypass:1}
La carte écrit la trame par accès direct à la mémoire (DMA) dans un tampon désigné par le prochain descripteur de réception et le marque comme terminé~; elle lève une interruption (peut-être
après le regroupement)~; le noyau programme le traitement NAPI dans une interruption logicielle, qui lit l'anneau, construit le tampon de socket, le fait passer
vers IP et UDP, le met en file sur le socket et réveille le lecteur~; \texttt{recv} le copie dans le tampon du programme.
\iqlookfor{DMA, l'anneau, l'interruption et le softirq, la file du socket, la copie et le réveil.}
\end{solution}

\begin{solution}{iq:nw:network-cards-and-kernel-bypass:2}
La carte retarde l'interruption de réception d'un délai ou d'un nombre de paquets, de sorte qu'une seule interruption serve plusieurs paquets~: moins de CPU, plus de latence.
Sur un hôte de données de marché, les anneaux du chemin critique sont scrutés, donc leur regroupement n'a pas d'importance~; pour les autres, choisir des valeurs modérées et
désactiver le regroupement adaptatif là où la prévisibilité compte.
\iqlookfor{le compromis latence-CPU, et le fait que le chemin critique l'évite par la scrutation.}
\end{solution}

\begin{solution}{iq:nw:network-cards-and-kernel-bypass:3}
Faire fonctionner le chemin de données en espace utilisateur~: accès brut à l'anneau (interfaces des fournisseurs, DPDK), \omterm{def:nw:network-cards-and-kernel-bypass:pmd}{piles réseau en espace utilisateur} derrière l'API des sockets, ou AF\_XDP.
Elle supprime les interruptions, les appels système, les copies et les réveils. On renonce aux outils du noyau, à l'isolation et au TCP (sauf si la pile apporte
le sien), et l'on consacre un cœur à chaque fil de scrutation.
\iqlookfor{les familles et les coûts, et pas seulement la vitesse.}
\end{solution}

\begin{solution}{iq:nw:network-cards-and-kernel-bypass:4}
Si le cœur de scrutation est vraiment isolé (autres tâches, interruptions, ticks de minuterie, threads du noyau qui s'y trouvent), la fréquence et les états de veille (C-state),
les défauts de page et les allocations sur le chemin, la taille de l'anneau face à la plus longue pause, et si la queue coïncide avec les rafales. Mesurer
avec des horodatages matériels, face aux horodatages logiciels.
\iqlookfor{d'abord les interférences d'ordonnancement, puis l'anneau.}
\end{solution}

\begin{solution}{iq:nw:network-cards-and-kernel-bypass:5}
La charge moyenne ne dit rien des pauses~: tant que l'hôte ne traite pas (une rafale, un désordonnancement, un réapprovisionnement par lots), la carte consomme
les descripteurs libres au rythme d'arrivée, et une pause plus longue que le nombre de descripteurs libres divisé par ce rythme fait perdre des paquets. Dimensionner l'anneau pour le
rythme des rafales multiplié par la plus longue pause, plus le lot de réapprovisionnement.
\iqlookfor{durée de pause fois rythme, et les descripteurs immobilisés par le traitement par lots.}
\end{solution}

\begin{solution}{iq:nw:network-cards-and-kernel-bypass:6}
L'horodatage du noyau est pris après la carte, son anneau, l'interruption et le pilote~; celui de la carte est pris au port. Leur
écart correspond au chemin de réception que la société cherche à régler, et seul l'horodatage de la carte se compare aux horodatages des commutateurs et
taps (chapitre~5).
\iqlookfor{où sur le chemin chaque horodatage est pris.}
\end{solution}
